% §3 Problem definition.  One continuous argument: the STDR recursion, the Theta(m^2) cost of its
% top-level call, the Bernoulli(p) observation model that answers that cost, and the entry-wise
% criterion eq:entrywise that every bound in the paper is measured against.  No subsections --
% the section is one argument and headings would restart it.
\section{Problem definition}\label{sec:problem}

The top-level call of \Cref{alg:stdr} governs the computational cost of the algorithm.
Computing the full matrix at the top level has a cost of order $\Theta(m^2\ell)$. For large trees, say, with
$m \sim 10^5$ terminal nodes and $\ell = 10^4$, the computational cost is on the order of $10^{14}$. 
In this paper, our goal is to derive a divide-and-conquer algorithm, with theoretical guarantees, with a computational cost that is linear in the number of terminal nodes, up to logarithmic terms. 

We apply a randomized approach, where only a random subset of the entries of the similarity matrix is computed at each level. 
Let $\Omega\in\{0,1\}^{m\times m}$ be a symmetric mask whose upper-triangular entries are
independent $\mathrm{Bernoulli}(p)$, $p\in(0,1]$. We compute a sparse similarity with
inverse-probability weighting \cite{candes2009matrix},
\begin{equation}\label{eq:sampling_estimator}
   \hat S_{ij}=\begin{cases}\tfrac1p\,\Omega_{ij}\,S_{ij}, & i\ne j,\\ 1, & i=j.\end{cases}.
\end{equation}
In \Cref{sec:assumptions} we prove, under appropriate assumptions, that a sample rate of $p=\Theta(\log m/m)$ suffices to reproduce the full-data split (\Cref{thm:main-sim}) with high probability. This reduces the overall computational cost of the divide-and-conquer algorithm to 
$\Theta(m\log m)$ pairwise estimates.


Our analysis addresses the asymptotic regime of infinite sequence length $\ell \to \infty$, which implies that the estimation error in the entries of $S$ reduces to $0$.
In this regime, \Cref{thm:stdr-split} guarantees that partitioning the terminal nodes according to the sign pattern of the Fiedler vector $v^{(2)}$ gives a valid partition $\Pi(S)$ in the underlying tree $\mathcal T$.
Our goal is to show that partitioning the terminal nodes according to the entry-wise sign pattern of the Fiedler vector $\hat v^{(2)}$, gives an identical partition with high probability. 
A sufficient condition is that the $\ell_{\infty}$ error between the two Fiedler vectors is bounded by 
\begin{equation}\label{eq:entrywise}
   \norm{\vvh-\vv}_\infty<\min_i|\vv_i|.
\end{equation}


This estimator is unbiased, $\E[\hat S]=S$, so $\EL=\hat L-L$ is mean zero, where
$\hat L=\hat D-\hat S$. The resulting error is missingness rather than additive noise:
most pairs are left at zero and the rest rescaled by $1/p$, so it is heavy, structured, and
its per-node variance grows as $p\to0$ (\Cref{prop:C-var}).

%\medskip
%\noindent\textbf{Problem.} Given the sparsely observed estimate $\hat S$ of
%\eqref{eq:sampling_estimator}, recover the Fiedler split $\Pisp{S}$ of the full matrix from the sign
%pattern of the Fiedler vector $\vvh$ of $\hat L$. We ask for the minimal sampling rate $p$ at which
%$\Pisp{\hat S}=\Pisp{S}$ \hp. Since $\E[\hat S]=S$, the question reduces to an
%eigenvector-perturbation problem over $\EL=\hat L-L$, which the remainder of the paper analyzes.
%The full-data split $\Pisp{S}$ is already a split of $\mathcal{T}$ (\Cref{thm:stdr-split}), so a
%rate at which $\Pisp{\hat S}=\Pisp{S}$ is a rate at which the sub-sampled operator returns a genuine
%split of the tree.


